\section{The Liquidity Provider's Decision Problem}
\label{sec:formulation}

We model the LP's problem as a risk-sensitive Markov decision process with
impulse control. The formalism follows from the economics. A rebalance is a
discrete intervention that costs gas and slippage rather than a continuous
control, and the volatility regime and order flow that drive the payoff are
latent. Decisions are taken every ten minutes over an episode of thirty days,
and the agent maps a fixed observation to a choice between holding and acting.
The reward measures the strategy's value against the fixed opening basket held
passively, net of the costs of the liquidity and hedge legs. Because the regime
and the flow are never observed directly, the process is partially observed,
and the observation supplies the history on which a decision can condition.
Figure~\ref{fig:loop} shows one decision cycle.

\begin{figure}[tbp]
\centering
\begin{tikzpicture}[
  font=\small,
  box/.style={draw, rounded corners=3pt, align=center, inner sep=5pt,
              fill=black!4, text width=3.1cm},
  arr/.style={-{Stealth[length=2.4mm]}, thick},
  lbl/.style={font=\scriptsize\itshape, midway, fill=white, inner sep=1.5pt}
]
  \node[box] (policy) at (0,2.1) {policy $\pi_\theta(a_t \mid o_t)$};
  \node[box] (action) at (5.1,2.1) {action $a_t$\\timing, width, hedge};
  \node[box] (market) at (5.1,0) {market step\\ten minutes, episode of thirty days};
  \node[box] (obs)    at (0,0)   {observation $o_{t+1}$\\24 features from past data};
  \draw[arr] (policy) -- node[lbl] {act} (action);
  \draw[arr] (action) -- node[lbl] {transition} (market);
  \draw[arr] (market) -- node[lbl] {reward $r_t$} (obs);
  \draw[arr] (obs) -- node[lbl] {observe} (policy);
\end{tikzpicture}
\caption{One cycle of the LP decision process. The policy acts every ten
minutes; a decision to rebalance moves the band and charges gas and slippage,
while a decision to hold leaves the position where it is. Rewards accumulate
over an episode of thirty days against the fixed opening basket.}
\label{fig:loop}
\end{figure}

\subsection{State}

The observation contains 24 features, computed from information available at
the decision time. Table~\ref{tab:observation} groups them by source. Band
geometry belongs in the observation because the action space is written in
tick spacings relative to the current tick. The distance to an edge sets the
time over which the band can earn fees, so a policy that cannot see how close
the price sits to an edge cannot time a rebalance. Heavy-tailed features such
as price, liquidity and gas are compressed with a signed logarithm, which
brings the dimensions onto a common scale for the network.

\begin{table}[tbp]
\caption{The 24 observation features, grouped by source. Every field is
computed from information available at the decision time, so an observation at
time $t$ cannot contain a bar from $t+1$.}
\label{tab:observation}
\centering
\small
\begin{tabular}{@{}L{0.18\linewidth}L{0.65\linewidth}r@{}}
\toprule
\textbf{Group} & \textbf{Features} & \textbf{n} \\
\midrule
Market & price, sqrt price, tick, 1\,h and 24\,h log returns, realised 24\,h volatility & 6 \\
Position & in-range flag, position value, two uncollected fee balances, IL against the fixed basket & 5 \\
Pool & active liquidity, fee tier & 2 \\
Costs & recent gas price, recent ETH reference price & 2 \\
Regime & regime label (bull, bear, sideways, high volatility) & 1 \\
Time & step in episode, steps remaining & 2 \\
Band & distance to each band edge and band width, in tick spacings & 3 \\
Hedge & short ratio, hedge equity, recent funding rate & 3 \\
\midrule
 & \textbf{Total} & \textbf{24} \\
\bottomrule
\end{tabular}
\end{table}

\subsection{Action}

At each decision the agent selects a factored action
\begin{equation}
a_t = (\kappa_t,\, c_t,\, w_t,\, h_t),
\label{eq:action}
\end{equation}
whose four heads are described in Table~\ref{tab:action}. Factoring the
action keeps the economics of the levers visible in the representation. The
study learns the width and hedge heads, while the timing head is exercised
once, at deployment, and the centre stays at the current tick. The restriction
follows from what the first experiments measured. On this pool the cost of
moving a band exceeds the fee that a move can recover, so the timing decision
collapses to the choice of whether to deploy at all. Section~\ref{sec:system}
documents the diagnostic runs behind the restriction, and
Section~\ref{sec:results} evaluates the learned policy under it.

The width ladder used for the study holds three arms that bracket the
economic trade-off. The narrow arm sits near the typical daily move, where fee
density is highest and the price is most likely to reach an edge. The middle
arm is just wider than the widest hand-written band. The wide arm approaches
the full range, where fee density is low but the band survives long stretches
of price travel. Together with the hold option these give the liquidity head
four choices, and the hedge head multiplies them by five ratios, so a full action
is one of $4\times5$ pairs. The interface also admits the wider timing and
centre grids, which remain available for the workstream on active timing.

\begin{table}[tbp]
\caption{The four action heads. Widths are given as half the band width, in
tick spacings. The general grid holds ten widths; the study's configuration
restricts the width head to three arms and the timing head to a single
deployment.}
\label{tab:action}
\centering
\small
\begin{tabular}{@{}L{0.14\linewidth}cL{0.28\linewidth}L{0.41\linewidth}@{}}
\toprule
\textbf{Head} & \textbf{Symbol} & \textbf{Options} & \textbf{Meaning} \\
\midrule
Rebalance & $\kappa_t$ & hold, act & spend gas and slippage to move the band, or leave it \\
Centre & $c_t$ & $-4,\dots,4$ spacings & shift the centre of the new band around the current tick \\
Width & $w_t$ & $\{1,2,5,10,25,50,$\newline $100,250,500,1000\}$ spacings & half of the new band's width \\
Hedge & $h_t$ & $\{0,0.25,0.5,0.75,1.0\}$ & share of the LP's ETH delta that the short offsets \\
\bottomrule
\end{tabular}
\end{table}

\subsection{Reward}

A rebalance converts inventory and pays for the privilege, so the reward
charges the move and credits the fees the new band earns. Writing $V_t$ for
the position value, $H_t$ for the fixed value of the opening deposit held
passively, and the hedge cash flows as mark PnL $\pi_t$, funding $\phi_t$ and
perp cost $\tau_t\le 0$, the reward is
\begin{equation}
r_t = \frac{1}{W_0}\Big(\mathrm{fees}_t - \mathrm{gas}_t - \mathrm{slip}_t
  + \Delta(V_t - H_t) + \pi_t + \phi_t + \tau_t - \lambda\,\hat{c}_t\Big),
\label{eq:reward}
\end{equation}
whose terms are defined in Table~\ref{tab:reward}. Dividing by the initial
capital $W_0$ keeps value targets at a workable scale without hiding any cost.
The benchmark is the \emph{fixed} opening basket. Because $H_t$ is never
rebased at a rebalance, the episode sum equals the excess over holding that
basket, which is the quantity the evaluation reports. Measuring against a
rebased entry price would instead credit the agent for value that the reported
statistic does not count.

\begin{table}[tbp]
\caption{The terms of the reward, in token0 (USDC) units before the $1/W_0$
scaling. Gas and slippage enter as positive magnitudes and carry a minus sign
in the reward.}
\label{tab:reward}
\centering
\small
\begin{tabular}{@{}L{0.16\linewidth}L{0.62\linewidth}c@{}}
\toprule
\textbf{Term} & \textbf{Meaning} & \textbf{Sign} \\
\midrule
$\mathrm{fees}_t$ & fee income earned in the step & $+$ \\
$\mathrm{gas}_t$ & gas for mints, burns, collects and rebalancing swaps & $-$ \\
$\mathrm{slip}_t$ & slippage on the inventory actually swapped & $-$ \\
$\Delta(V_t-H_t)$ & change in position value against the fixed opening basket & $\pm$ \\
$\pi_t$ & mark PnL of the short perpetual & $\pm$ \\
$\phi_t$ & funding paid or received on the short & $\pm$ \\
$\tau_t$ & perp taker cost of adjusting the short & $\le 0$ \\
$\lambda\,\hat{c}_t$ & penalty on the trailing expected shortfall of per-step losses, $\lambda=0.05$ in the study & $-$ \\
\bottomrule
\end{tabular}
\end{table}

\subsection{Risk and benchmarks}

The learning objective is the discounted return with a penalty on the recent
tail of per-step losses. The risk term in Eq.~\eqref{eq:reward} subtracts
$\lambda$ times the recent $5\,\%$ expected shortfall of per-step equity
changes, measured over the last trading day, so deeper tail losses cost more
reward. The penalty is computed on the joint equity of the liquidity and
hedge legs, so the hedge ratio can move the quantity the policy is asked to
control. The default configuration disables the term ($\lambda=0$); the
study's configuration sets $\lambda=0.05$, small next to a full tail loss.
The evaluation reports the same quantity the penalty controls. With losses
$L=-G$ defined on the episode return $G$, the conditional value-at-risk at
level $\alpha$ is
\begin{equation}
\mathrm{CVaR}_\alpha(L) = \min_{\nu\in\mathbb{R}}
\left\{\nu+\frac{1}{1-\alpha}\mathbb{E}\big[(L-\nu)^+\big]\right\},
\label{eq:cvar}
\end{equation}
the Rockafellar--Uryasev representation\cite{rockafellar2000cvar}, which we
report at $\alpha=0.95$. Evaluation therefore reports the annualised Sharpe
ratio and CVaR at $5\,\%$ as primary statistics, with raw PnL and excess over
HODL as secondary ones. Comparisons use paired per-window differences and
confidence intervals across seeds.

The hedge is a short perpetual position sized against the LP's ETH exposure.
Its mark PnL over a step is the negative product of the short notional and the
price change, plus funding and the taker cost of adjusting the position. Three
reference points make the hedge leg interpretable. The unhedged LP carries the
full directional exposure, the policy chooses a hedge ratio step by step, and
a hedged HODL portfolio holds a delta-matched short. The third is the
benchmark that prevents a claim of success against a weaker opponent. A policy
that only removes directional risk must still beat a passive portfolio with
the same risk profile. The overlay is enabled in the study's configuration,
and the hedge ratio is a learned decision rather than a hand-set constant.
